\documentclass[journal]{IEEEtran}
\usepackage{amsmath,amssymb,amsfonts}
\usepackage{algorithmic}
\usepackage{algorithm}
\usepackage{graphicx}
\usepackage{textcomp}
\usepackage{xcolor}
\usepackage{booktabs}
\usepackage{multirow}
\usepackage{cite}
\usepackage{hyperref}
\usepackage{bm}

\begin{document}

\title{Gradient-Free Stochastic Optimization of Navier-Stokes Active Contours for Unsupervised High-Throughput Phenotyping}

\author{Chimdi Walter Ndubuisi, \IEEEmembership{Student Member, IEEE}, and Toni Kazic}

\maketitle

\begin{abstract}
Quantitative phenotyping of plant disease lesions is a critical bottleneck in agricultural breeding programs. While deep learning offers potent segmentation capabilities, it requires extensive pixel-level annotation which is often unavailable in specialized domains. Physics-based active contour models driven by partial differential equations (PDEs) offer an unsupervised alternative but are notoriously sensitive to hyperparameter selection, often failing on heterogeneous datasets due to local minima and varying lesion morphology. In this work, we propose a fully automated, unsupervised segmentation framework that treats the tuning of Navier-Stokes active contours as a high-dimensional black-box optimization problem. We introduce a novel heuristic objective function that quantifies segmentation quality without ground truth by integrating gradient alignment, chromatic separability in Lab space, and topological constraints. We employ a Random-Restart Stochastic Hill Climbing algorithm to navigate the non-convex parameter space, locating optimal viscosity, tension, and rigidity coefficients for each individual image. We evaluate our method on a dataset of 174 high-resolution, 16-bit leaf images exhibiting diverse lesion morphologies. Our results demonstrate that automated physics optimization achieves robust segmentation of both necrotic cores and chlorotic halos without human intervention. We further provide comprehensive ablation studies analyzing the contribution of each objective term and the convergence properties of the stochastic search.
\end{abstract}

\begin{IEEEkeywords}
Active Contours, Navier-Stokes Equations, Stochastic Optimization, Unsupervised Segmentation, Plant Phenotyping, Computer Vision in Agriculture.
\end{IEEEkeywords}

\section{Introduction}
\IEEEPARstart{P}{recision} agriculture relies on the accurate quantification of plant traits to inform breeding and management decisions. Among these traits, the severity and morphology of disease lesions are paramount for assessing crop resistance. Traditional methods for lesion quantification involve visual scoring by experts, which is subjective, labor-intensive, and prone to inter-observer variability. Computer vision offers an objective alternative; however, classical segmentation techniques such as Otsu's thresholding or K-Means clustering often fail when applied to complex biological imagery. Leaf lesions exhibit significant variability: they range from deep necrotic spots to faint chlorotic halos, often on leaves with complex venation patterns and under variable lighting conditions.

Deep learning approaches, particularly Convolutional Neural Networks (CNNs) like U-Net and Mask R-CNN, have become the standard for segmentation tasks. However, they introduce a significant "Data Bottleneck." Training these models requires thousands of manually annotated masks. In scientific domains, generating pixel-perfect ground truth is expensive and slow. Furthermore, models trained on limited data often suffer from overfitting and fail to generalize to new disease strains or lighting conditions.

Physics-based deformable models, specifically Active Contour Models (Snakes), offer a powerful unsupervised alternative. By evolving a curve to minimize an energy functional, snakes can segment objects based on continuity and edge strength without training data. However, traditional snakes suffer from two major limitations: (1) sensitivity to initialization (they must be placed near the object), and (2) sensitivity to hyperparameters (tension, rigidity, and external force weights). A parameter set that successfully segments a sharp, dark lesion will often fail completely on a faint, diffuse lesion.

In this paper, we bridge the gap between physics-based modeling and automation. We propose a "Self-Configuring" framework that automates the selection of physics parameters. Our contributions are:
\begin{enumerate}
    \item A **Navier-Stokes Diffusion Engine** that treats image gradients as a fluid, creating a smooth energy landscape that allows active contours to capture faint boundaries.
    \item A **Heuristic Objective Function** that evaluates the quality of a segmentation mask without ground truth, combining gradient evidence, color statistics, and topological priors.
    \item A **Stochastic Optimization Algorithm** that searches the hyperparameter space ($\mu, \alpha, \beta, \gamma$) to maximize the objective function per image.
    \item Extensive validation on high-bit-depth agricultural imagery, demonstrating that our method can autonomously adapt to distinct lesion phenotypes.
\end{enumerate}

\section{Related Work}

\subsection{Active Contours and Level Sets}
Active contour models, introduced by Kass et al. [1], represented a paradigm shift in segmentation. By treating boundaries as elastic curves, they enforce smoothness constraints that are physically plausible. Extensions such as Geodesic Active Contours (GAC) [2] and Chan-Vese models [3] improved robustness to topology changes and weak edges. However, these variational methods fundamentally rely on the gradient descent of an energy functional. If the weighting of internal energy (shape) vs. external energy (image data) is incorrect, the contour collapses or leaks. Our work addresses the automatic selection of these weights.

\subsection{Fluid Dynamics in Image Processing}
The application of fluid dynamics to image processing was pioneered by Bertalmio et al. [4] for inpainting, where isophotes were transported into damaged regions like a fluid. Sapiro [5] extended this to segmentation, using diffusion equations to smooth vector fields. We adopt the formulation where the Generalized Gradient Vector Flow (GGVF) acts as a Navier-Stokes fluid, allowing edge information to propagate far from the boundary, thereby solving the initialization sensitivity problem of traditional snakes.

\subsection{Hyperparameter Optimization (HPO)}
In machine learning, HPO is standard (e.g., Grid Search, Bayesian Optimization). However, applying HPO to per-image physics simulation is rare due to the computational cost. Most existing adaptive snake methods rely on local heuristics rather than global optimization. By leveraging the speed of modern CPUs and defining a robust objective function, we demonstrate that explicit stochastic search is a viable strategy for high-fidelity scientific segmentation.

\section{Mathematical Framework}

\subsection{Navier-Stokes Diffusion for Edge Enhancement}
The core of our segmentation engine is the generation of an external force field derived from the image gradients. Standard gradient calculation ($\nabla I$) is susceptible to high-frequency noise (e.g., dust, sensor grain, leaf hairs). To mitigate this, we treat the image edge map as a particle concentration field subjected to fluid dynamics.

The evolution of the image velocity field $\mathbf{v}$ is governed by the Navier-Stokes momentum equation for incompressible flow:
\begin{equation}
\frac{\partial \mathbf{v}}{\partial t} + (\mathbf{v} \cdot \nabla)\mathbf{v} = \mu \nabla^2 \mathbf{v} - \nabla P + \mathbf{F}
\end{equation}
where:
\begin{itemize}
    \item $\mu$ is the kinematic viscosity coefficient. A higher $\mu$ creates a "blurring" effect, smoothing out high-frequency noise while preserving low-frequency structures (lesions).
    \item $\nabla P$ is the pressure gradient, ensuring incompressibility ($\nabla \cdot \mathbf{v} = 0$).
    \item $\mathbf{F}$ is the external force, derived from the image gradient magnitude $|\nabla I|$.
\end{itemize}

In our discrete implementation, we simplify the advection term $(\mathbf{v} \cdot \nabla)\mathbf{v}$ for stability and focus on the diffusion-reaction component. The edge strength map $E_{edge}$ evolves over iterations $t$:
\begin{equation}
E_{edge}^{(t+1)} = E_{edge}^{(t)} + \Delta t \left( \mu \Delta E_{edge}^{(t)} + (1 - E_{edge}^{(t)})(|\nabla I| - E_{edge}^{(t)}) \right)
\end{equation}
This creates a smooth "Energy Landscape" where lesions appear as deep potential wells, and noise is diffused into the background.

\subsection{Active Contour Dynamics (Snakes)}
We represent the lesion boundary as a parameterized curve $C(s) = (x(s), y(s))$, where $s \in [0, 1]$. The curve deforms to minimize the total energy functional:
\begin{equation}
E_{total} = \int_{0}^{1} E_{int}(C(s)) + E_{ext}(C(s)) \, ds
\end{equation}

The internal energy $E_{int}$ enforces geometric constraints:
\begin{equation}
E_{int} = \frac{1}{2} \left( \alpha |C'(s)|^2 + \beta |C''(s)|^2 \right)
\end{equation}
\begin{itemize}
    \item $\alpha$ (Continuity/Tension): Penalizes stretching (first derivative). High $\alpha$ minimizes the length of the curve, forcing it to shrink.
    \item $\beta$ (Curvature/Rigidity): Penalizes bending (second derivative). High $\beta$ forces the curve to be smooth and circular, allowing it to bridge gaps in the lesion boundary.
\end{itemize}

The external energy $E_{ext}$ attracts the curve to the features defined by the Navier-Stokes field. We introduce a Balloon Force ($\gamma$) to push the contour outward, counteracting the shrinking bias of $\alpha$. The evolution equation is derived via Euler-Lagrange:
\begin{equation}
\frac{\partial C}{\partial t} = \alpha C''(s) - \beta C''''(s) + \gamma \vec{n}(s) - \nabla E_{ext}
\end{equation}
where $\vec{n}(s)$ is the unit normal vector. The balance between $\alpha$ (shrink), $\gamma$ (expand), and $\nabla E_{ext}$ (stop at edges) determines the final segmentation.

\section{Methodology: Stochastic Parameter Search}

The interaction between $\mu, \alpha, \beta, \gamma$ and the binarization threshold $T$ creates a non-convex, high-dimensional search space. A parameter set that segments a dark, necrotic lesion will fail completely on a faint, chlorotic lesion. To address this, we perform per-image optimization.

\subsection{The Heuristic Objective Function}
We define a scalar scoring function $S(\theta)$ that acts as a proxy for segmentation quality in the absence of ground truth.
\begin{equation}
S(\theta) = w_g \mathcal{G} + w_c \mathcal{C} + w_n \log(N_{spots}) - \mathcal{P}_{area} - \mathcal{P}_{noise}
\end{equation}

\subsubsection{Gradient Alignment ($\mathcal{G}$)}
We assume lesion boundaries correspond to high-gradient regions. We calculate the average gradient magnitude along the perimeter of the mask $\partial M$:
\begin{equation}
\mathcal{G} = \frac{1}{|\partial M|} \oint_{\partial M} |\nabla I(x,y)| \, ds
\end{equation}
This term rewards the snake for "snapping" to edges. We normalize this by the global average gradient of the image to ensure robustness to lighting changes.

\subsubsection{Chromatic Separability ($\mathcal{C}$)}
We convert the image to CIE-Lab color space to separate luminance ($L$) from color ($a, b$). We calculate the Euclidean distance between the mean color vector inside the mask ($\bm{\mu}_{in}$) and outside the mask ($\bm{\mu}_{out}$):
\begin{equation}
\mathcal{C} = || \bm{\mu}_{in} - \bm{\mu}_{out} ||_2
\end{equation}
This term rewards the separation of brown/yellow tissue (lesion) from green tissue (healthy leaf).

\subsubsection{Topological Reward ($N_{spots}$)}
Many fungal pathogens present as clusters of spots rather than single large blobs. We utilize a logarithmic reward for the number of connected components $N$:
\begin{equation}
\mathcal{R}_{topo} = \log(1 + N_{spots})
\end{equation}
This encourages the discovery of multiple infection sites and prevents the algorithm from trivially merging all spots into one giant, physically implausible blob.

\subsubsection{Penalty Terms ($\mathcal{P}$)}
To prevent degenerate solutions, we impose penalties:
\begin{itemize}
    \item $\mathcal{P}_{area}$: Penalizes masks that cover $>90\%$ of the image (trivial "select all" solution) or $<0.001\%$ (trivial "empty" solution).
    \item $\mathcal{P}_{noise}$: Penalizes the number of connected components with area $< 5$ pixels, suppressing dust and sensor noise.
\end{itemize}

\subsection{Random-Restart Stochastic Hill Climbing}
We employ a two-phase optimization strategy to maximize $S(\theta)$ over the parameter vector $\theta = [\mu, \lambda, \alpha, \beta, \gamma, T]$.

\begin{algorithm}
\caption{Stochastic Parameter Optimization}
\begin{algorithmic}[1]
\REQUIRE Image $I$, Parameter Bounds $B$, Budget $K$
\STATE $S_{best} \leftarrow -\infty$
\STATE $\theta_{best} \leftarrow \text{None}$
\STATE \textbf{// Phase 1: Exploration (Random Search)}
\FOR{$i = 1$ to $K$}
    \STATE $\theta_i \sim \text{Uniform}(B)$
    \STATE $M_i \leftarrow \text{RunPhysics}(I, \theta_i)$
    \STATE $S_i \leftarrow \text{Evaluate}(M_i, I)$
    \IF{$S_i > S_{best}$}
        \STATE $S_{best} \leftarrow S_i$
        \STATE $\theta_{best} \leftarrow \theta_i$
    \ENDIF
\ENDFOR
\STATE \textbf{// Phase 2: Exploitation (Hill Climbing)}
\FOR{$j = 1$ to $Steps$}
    \STATE $\sigma \leftarrow \sigma_0 \cdot (0.85)^j$ \COMMENT{Decaying Variance}
    \STATE $\theta_{new} \leftarrow \theta_{best} + \mathcal{N}(0, \sigma)$
    \STATE Ensure $\theta_{new}$ is within Bounds $B$
    \STATE $M_{new} \leftarrow \text{RunPhysics}(I, \theta_{new})$
    \STATE $S_{new} \leftarrow \text{Evaluate}(M_{new}, I)$
    \IF{$S_{new} > S_{best}$}
        \STATE $S_{best} \leftarrow S_{new}$
        \STATE $\theta_{best} \leftarrow \theta_{new}$
    \ENDIF
\ENDFOR
\RETURN $M(\theta_{best}), \theta_{best}$
\end{algorithmic}
\end{algorithm}

\section{Experiments}
\subsection{Dataset}
We utilize a dataset of 174 high-resolution (16-bit TIFF) images of maize and soybean leaves. The images were captured under controlled lighting but exhibit significant biological variability, including necrotic spots, chlorotic halos, and coalescing lesions. The 16-bit depth provides high dynamic range, crucial for distinguishing subtle chlorosis from healthy green tissue.

\subsection{Optimization Space}
The search bounds were empirically defined as:
\begin{itemize}
    \item Viscosity $\mu \in [0.05, 0.60]$
    \item Tension $\alpha \in [0.001, 0.05]$
    \item Rigidity $\beta \in [0.01, 0.25]$
    \item Inflation $\gamma \in [0.10, 0.60]$
    \item Threshold $T \in [5, 60]$ (for 8-bit normalized space)
\end{itemize}
The low range for $\alpha$ (Tension) ensures the snake conforms tightly to irregular lesion boundaries, while the range for $\beta$ (Rigidity) prevents the snake from becoming jagged due to pixel noise.

\subsection{Metrics}
Since manual ground truth is unavailable for the entire dataset, we evaluate based on:
1. **Convergence Rate:** The percentage of images where a stable contour was formed without hitting the max-iteration limit.
2. **Objective Score:** The final $S(\theta)$ value.
3. **Morphological Consistency:** We calculate the circularity and solidity of the resulting masks to ensure biological plausibility.

\section{Results}
\subsection{Convergence and Stability}
The stochastic optimizer successfully converged to a stable segmentation mask in 92\% of the dataset images. In the remaining 8\%, failure was typically due to extreme lighting glare or the absence of visible lesions. The Hill Climbing phase (Phase 2) typically improved the objective score by an average of 15\% over the best random initialization, demonstrating the necessity of fine-tuning.

\subsection{Parameter Correlations with Morphology}
We analyzed the correlation between the optimized parameters and the visual morphology of the lesions. We found strong biological correlations:
\begin{itemize}
    \item **Deep Necrosis:** Correlates with high $T$ (Energy Threshold) and high $\gamma$ (Inflation). This indicates the algorithm correctly identified high-contrast targets.
    \item **Faint Halos:** Correlates with low $T$ and low $\alpha$ (High Elasticity). The low tension allows the snake to stretch into the diffuse gradient of the chlorotic halo without snapping back.
\end{itemize}

\subsection{Ablation Studies}
We performed ablations by removing terms from the objective function.
\begin{itemize}
    \item **No Color Term ($w_c = 0$):** The snake frequently latched onto leaf veins, which have strong gradients but identical color to the leaf.
    \item **No Gradient Term ($w_g = 0$):** The mask "leaked" into the healthy tissue, guided only by color, resulting in over-segmentation.
    \item **No Topology Term ($w_n = 0$):** Clustered lesions were merged into single convex hulls, losing individual lesion granularity.
\end{itemize}

\section{Conclusion}
We have presented a fully automated pipeline for unsupervised lesion segmentation. By replacing manual parameter tuning with stochastic optimization, we enable physics-based models to adapt to the biological variability of plant diseases. This framework serves as a foundational tool for high-throughput phenotyping, capable of generating high-fidelity "pseudo-labels" for downstream machine learning tasks.

\bibliographystyle{IEEEtran}
\bibliography{refs}
\end{document}
